\bookchapter{Capitolo 32: Manuale di pronto soccorso per la sicurezza del telefono}

\chapterintro{Questo capitolo è un «piano di emergenza». Speriamo non ti serva mai, ma se perdi il telefono, te lo rubano o subisci una truffa, questo manuale ti aiuta a ridurre i danni al minimo; in caso di terremoti, pioggia intensa o altre calamità, il telefono è anche lo strumento per ricevere allerte, localizzarti e restare in contatto. Ti consigliamo di leggerlo adesso, poi fai uno screenshot dei passaggi chiave o stampalo e tienilo in casa: quando succede qualcosa il telefono potrebbe non essere più con te, e in caso di calamità potrebbe essere scarico o senza segnale.}

\booksection{Scenario zero: terremoto / maltempo estremo — come il telefono può salvarti la vita}

Qui si protegge la sicurezza personale, non il denaro sugli account. Per le impostazioni di allerta vedi Capitolo 13 «Allerte meteo, calamità naturali e terremoti»; qui parliamo solo di cosa fare dopo lo scatto dell'allarme.

\textbf{Scatta l'allarme terremoti: prima riparati, poi evacua, non pensare solo a filmare}\par

\begin{enumerate}
  \item Riparati subito: buttati a terra, copriti, aggrappati (sotto il tavolo, angoli di muri portanti, mobili solidi), proteggi testa e collo; tieniti lontano da vetri, pareti esterne e lampadari.
  \item Non buttarti sulle scale: la scossa principale dura spesso solo decine di secondi e le scale possono essere più pericolose; l'ascensore è assolutamente vietato.
  \item Evacua in ordine quando la scossa finisce: prendi le scale, attento alle repliche e agli oggetti che cadono; vai in spazio aperto, lontano da pali, cartelloni ed edifici pericolosi.
  \item Usa il telefono con intelligenza: un messaggio breve ai contatti di emergenza («Sono a XX, al sicuro/ferito, ho bisogno di aiuto») risparmia più batteria di una lunga chiamata; se puoi, preferisci testo/SMS/WeChat.
  \item Senza segnale: vedi Capitolo 7 SOS satellitare / messaggi satellitari (se il modello lo supporta); oppure prova in alto o in spazio aperto; in alcuni casi è utile anche il walkie-talkie Bluetooth senza rete.
  \item Disattiva le funzioni non necessarie: in emergenza attiva la modalità aereo, poi aprila a orari programmati, spegni la scansione Bluetooth e abbassa la luminosità per resistere qualche ora in più; riserva l'ultimo bit di batteria per dire che stai bene e per i messaggi di soccorso.
\end{enumerate}

\begin{tcbwarn}{Pioggia intensa / tifone / colate detritiche: agisci già in fase di allerta}
\begin{itemize}
  \item Con allerta rossa: riduci al minimo le uscite, tieniti lontano da spazi interrati, sottopassi, fiumi e pendii.
  \item Se sei bloccato all'aperto: dirigiti verso l'alto di edifici solidi, non attraversare a piedi l'acqua (potrebbero mancare tombini o esserci dispersioni elettriche); chiamando 119/110 indica punti di riferimento e livello dell'acqua.
  \item Prima di partire: controlla i rischi con l'app meteo e gli account di allerta locali; nei lunghi viaggi in auto evita le strade di montagna lungo fiumi e dirupi.
  \item Accordo in famiglia: stabilite «se perdiamo il contatto, facciamo fede all'ultima posizione su WeChat/SMS», per evitare di esaurire la batteria chiamando ossessivamente.
\end{itemize}

\end{tcbwarn}

\begin{tcbpractice}{4 cose da fare già da oggi}
\begin{enumerate}
  \item Attiva l'allerta terremoti di sistema o un'app di allerta affidabile e consentile di superare il silenzio.
  \item Metti in evidenza 2-3 «contatti di emergenza» in rubrica, visibili anche dalla schermata di blocco (scheda medica iOS/Android o informazioni di blocco).
  \item Tieni un power bank in casa e nella borsa; prima dei lunghi viaggi caricalo e porta il cavo.
  \item Memorizza almeno un numero da poter comporre senza rubrica (il cellulare di un familiare) e scrivilo su carta nel portafoglio.
\end{enumerate}

\end{tcbpractice}

\booksection{Scenario uno: telefono perso/rubato — i 30 minuti d'oro}

I primi 30 minuti dopo la perdita sono «tempo d'oro»: se agisci bene proteggi i tuoi soldi e la tua privacy. Segui questo ordine senza saltare passaggi:

\textbf{Procedura standard nei 30 minuti dopo la perdita}\par

\begin{enumerate}
  \item Chiama il tuo numero da un altro telefono (minuto 1): potrebbe essere caduto lì vicino e una persona onesta può restituitelo.
  \item Blocca da remoto il telefono (minuti 2-5): da un altro telefono o computer accedi al servizio di localizzazione: iPhone: iCloud.com > Trova > seleziona il tuo dispositivo > Segna come smarrito
  \item Android per marca: Huawei cloud.huawei.com / Xiaomi i.mi.com / OPPO cloud.oppo.com ecc.
  \item Sospendi la SIM (minuti 5-10): chiama l'assistenza dell'operatore (China Mobile 10086 / Unicom 10010 / Telecom 10000) e chiedi di sospendere il servizio SIM. È il passaggio più critico: se i truffatori hanno la tua SIM possono ricevere i codici di verifica e reimpostare tutte le password.
  \item Congela i conti di pagamento (minuti 10-20): Alipay: chiama 95188 per il blocco, oppure da un altro telefono accedi ad Alipay > Il mio > Impostazioni > Centro sicurezza > Congelamento di emergenza
  \item WeChat: dopo la denuncia al 110 contatta l'assistenza WeChat, oppure Centro sicurezza WeChat > Blocca account
  \item App bancarie: chiama il numero sul retro della carta e chiedi di congelare il home banking
  \item Avvisa amici e parenti (minuti 20-25): pubblica su Moments o fai pubblicare un messaggio in famiglia per dire che hai perso il telefono, così i truffatori non possono spacciarsi per te.
  \item Denuncia (minuti 25-30): chiama 110 o vai al commissariato più vicino e ottieni la ricevuta di denuncia. Servirà per assicurazioni e controversie bancarie.
\end{enumerate}

\booksection{Scenario due: sei stato truffato — 72 ore d'oro per agire}

Se ti rendi conto di essere stato truffato, la cosa più importante è restare calmo. Decisioni prese nel panico causano danni doppi. Segui questi passaggi:

\begin{enumerate}
  \item Smetti immediatamente di trasferire denaro: qualunque cosa dica il truffatore, non inviare un altro centesimo. «Invia un altro bonifico di garanzia e ti restituiamo i soldi di prima» è la classica truffa a catena.
  \item Conserva tutte le prove: non cancellare le chat! Non cancellare le registrazioni delle chiamate! Fai screenshot di chat, bonifici e dati dell'account dell'altra parte. Sono prove fondamentali per la denuncia e il recupero del denaro.
  \item Chiama subito il 96110: è la linea anti-frode nazionale. L'operatore ti guiderà nel blocco di emergenza del pagamento: se il bonifico è recente (di solito entro 24 ore) la banca può congelare il conto destinatario e intercettare i fondi.
  \item Chiama 110 per denunciare: porta telefono, documento e ricevute del bonifico al commissariato più vicino per la deposizione.
  \item Contatta banca/piattaforma di pagamento: comunica la truffa e chiedi di congelare i conti coinvolti. Se è una carta di credito clonata, chiedi alla banca di avviare la procedura di chargeback.
  \item Cambia tutte le password: da un altro dispositivo sicuro cambia le password di Apple ID/account del marchio, WeChat, Alipay e app bancarie e attiva l'autenticazione a due fattori.
\end{enumerate}

\booksection{Scenario tre: il telefono ha un virus/trojan — cosa fare?}

Se noti questi sintomi — annunci che spuntano dal nulla, app strane che si installano da sole, costi in bolletta alle stelle, consumo di dati anomalo — è molto probabile che ci sia malware:

\begin{itemize}
  \item Passo 1: modalità provvisoria — su Android tieni premuto il tasto di accensione per spegnere, poi tienilo premuto per riaccendere; quando vibra, tieni premuto il tasto volume giù finché non compare la «Modalità provvisoria» (in basso a sinistra). In modalità provvisoria tutte le app di terze parti sono disattivate e puoi individuare il problema.
  \item Passo 2: disinstalla le app sospette — in modalità provvisoria vai in Impostazioni > Gestione app, ordina per data e disinstalla le app sconosciute installate di recente.
  \item Passo 3: installa un'app di sicurezza ed esegui una scansione — scarica «Tencent Mobile Manager» o «360 Mobile Guard» e fai una scansione completa.
  \item Passo 4: ripristino di fabbrica (soluzione estrema) — se i metodi precedenti non funzionano, esegui il backup dei dati importanti e ripristina le impostazioni di fabbrica. Attenzione: il ripristino cancella tutti i dati, salva prima foto, rubrica e altri file importanti.
  \item iPhone quasi non si infetta — il meccanismo sandbox di iOS rende difficile la sopravvivenza ai malware. Se il tuo iPhone si comporta in modo strano, quasi sempre è un bug di un'app o un problema di impostazioni, non un virus. Il ripristino di fabbrica risolve il 99\% dei problemi.
\end{itemize}

\booksection{Scenario quattro: account rubato — riconquista rapidamente il controllo}

\begin{enumerate}
  \item Cambia subito la password: se riesci ancora ad accedere, cambiala immediatamente. La nuova password deve essere completamente diversa da quella vecchia.
  \item Disconnetti gli altri dispositivi: nelle impostazioni di sicurezza dell'account trova «Gestione dispositivi di accesso» e disconnetti tutti i dispositivi sconosciuti.
  \item Controlla i registri di accesso: esamina gli accessi degli ultimi 7 giorni; se ci sono accessi da luoghi strani, l'account è effettivamente compromesso.
  \item Attiva l'autenticazione a due fattori: se non l'hai ancora attivata, fallo subito. È il metodo più efficace per evitare nuovi furti di account.
  \item Controlla i dati collegati: verifica se il numero di telefono o l'email collegati all'account sono stati modificati. In tal caso, recupera l'account con la funzione di ricorso.
  \item Controlla le impostazioni di pagamento: se sono stati rubati WeChat/Alipay, verifica che i pagamenti senza password e gli addebiti automatici non siano stati manomessi e che non siano state aggiunte carte a conti sconosciuti.
\end{enumerate}

\booksection{Scenario cinque: hai ricevuto un messaggio di estorsione — non avere paura, non pagare}

Le estorsioni più comuni sono due: (1) «Ho un tuo video/foto intimo, se non paghi lo mando a tutti i tuoi contatti»; (2) «Il tuo telefono è stato hackerato, se non pago lo blocco». In questi casi:

\begin{itemize}
  \item Al 99,9\% è una truffa: non hanno alcun dato sensibile tuo. Inviare lo stesso messaggio di estorsione a decine di migliaia di persone, qualcuno pagherà per paura.
  \item Non rispondere, non pagare: rispondere significa dire «questo numero è attivo» e diventerai un «bersaglio estorribile» con ancora più disturbi.
  \item Fai uno screenshot e conserva le prove: poi blocca il mittente.
  \item Cambia le password: se sei preoccupato, cambia la password dei tuoi account principali per stare più tranquillo.
  \item Se c'è davvero il rischio di una fuga di dati sensibili: denuncia. La polizia ha esperienza in questi casi.
\end{itemize}

\booksection{Scenario sei: truffe con deepfake vocali e video AI — ciò che vedi non è più vero}

Dal 2024 le tecnologie di deepfake e sintesi vocale AI permettono ai truffatori di «diventare» tuo figlio, un parente o un superiore in videochiamata. «Mamma, ho avuto un incidente, mandami subito i soldi» — la persona nel video sembra e suona esattamente tuo figlio, ma è falsa. È una delle truffe più difficili da individuare e più pericolose.

\textbf{Passo 1: riconoscere le falle — quattro punti deboli comuni del deepfake}\par

\begin{itemize}
  \item Frequenza di battito delle palpebre anomala: una persona reale sbatte le palpebre 15-20 volte al minuto; nei deepfake si sbatte troppo poco o in modo rigido.
  \item Labbra e audio non sincronizzati: tra il movimento della bocca e il suono c'è uno scarto di frazioni di secondo.
  \item Bordi del volto sfocati: il contorno del viso, le orecchie e i capelli possono presentare «frange» o lampeggi.
  \item Direzione della luce incoerente: i bagliori e le ombre sul volto non combaciano con la luce di sfondo, oppure di profilo i tratti «girano» in modo innaturale.
\end{itemize}

\begin{tcbinfo}{Passo 2: tre modi per verificare l'autenticità}
\begin{itemize}
  \item Chiedi un'azione specifica: fai girare la testa, coprire il viso o la bocca con la mano e poi aprire e chiudere. Il deepfake si tradisce con le occlusioni e i grandi angoli di rotazione: il viso si deforma o non segue il movimento.
  \item Fai una domanda che solo l'altra persona conosce: inventa una domanda falsa, ad esempio «quale clinica abbiamo portato il cane l'ultima volta?» — un truffatore non saprà rispondere o balbetta.
  \item Richiama un numero noto: chiudi la video e chiama direttamente il numero salvato in rubrica. Il truffatore non può controllare il vero cellulare: si scopre subito.
\end{itemize}

\end{tcbinfo}

\begin{tcbwarn}{Passo 3: cosa fare dopo aver confermato la truffa}
\begin{enumerate}
  \item Smetti subito di trasferire denaro: per quanto «urgente», fermati. Chi ti spinge a «bonificare subito e non riagganciare» è quasi sempre un truffatore.
  \item Conserva tutte le prove: registra lo schermo, fai screenshot di videochiamata, chat, account e bonifici — sono materiali fondamentali per la perizia tecnica.
  \item Chiama la linea anti-frode 96110: specifica che si tratta di truffa con deepfake vocale/video AI; l'operatore avvierà il blocco di emergenza.
  \item Chiama 110 per denunciare: porta tutte le prove al commissariato e sottolinea l'uso di tecnologie AI, così la polizia può classificare il reato.
  \item Avvisa amici e parenti: i truffatori potrebbero usare il tuo video falso per ingannare altri; pubblica un avviso su Moments o nel gruppo di famiglia.
\end{enumerate}

\end{tcbwarn}

\booksection{Scenario sette: attacco SIM swap (sequestro SIM) — il tuo numero è stato «rubato»}

È un attacco poco conosciuto ma devastante. Il SIM swap consiste nel fatto che l'aggressore, falsificando i tuoi dati, si spaccia per te in filiale e richiede una nuova SIM. Appena la SIM viene sostituita, il tuo telefono perde improvvisamente il segnale e il telefono dell'aggressore, con la nuova SIM, riceve tutti i tuoi codici di verifica: Alipay, banca, WeChat, tutto spalancato.

\begin{tcbinfo}{Come accorgersene? Tre segnali di allarme del SIM swap}
\begin{itemize}
  \item Il telefono mostra improvvisamente «Nessun servizio»/«Nessuna SIM» anche se non l'hai toccata e il riavvio non serve: spesso è l'istante del sequestro.
  \item Ricevi un SMS dell'operatore «Sostituzione SIM completata» o «Servizio attivato» — ma tu non hai fatto nulla.
  \item Improvvisamente non ricevi più i codici di verifica, e nemmeno altri SMS: i messaggi sono stati dirottati sulla SIM dell'aggressore.
\end{itemize}

\end{tcbinfo}

\begin{tcbwarn}{Misure di protezione — metti «sotto chiave» anche la SIM}
\begin{itemize}
  \item Imposta il PIN della SIM: percorso Impostazioni > Telefono > PIN SIM. Attivato, ogni riavvio o cambio SIM richiede il PIN: anche se qualcuno ha la carta, non la può usare.
  \item Attiva la verifica a due fattori dell'operatore: nell'app China Mobile/Unicom/Telecom trova «Centro sicurezza» o «Verifica servizi» e attiva l'avviso per sostituzioni SIM non richieste e la doppia password per le pratiche. Così in filiale verrà richiesta un'ulteriore verifica d'identità.
  \item Account importanti non solo con SMS: per banche, Apple ID e altri account di valore preferisci app autenticatrici (Google Authenticator, Microsoft Authenticator) o chiavi hardware invece dei codici SMS.
  \item Filigrana sulle fotocopie dei documenti: sulle copie per pratiche scrivi «solo per uso XX» — gli aggressori spesso si spacciano per te con documenti d'identità trapelati.
\end{itemize}

\end{tcbwarn}

\textbf{Gestione di emergenza dopo il sequestro — la velocità è tutto}\par

\begin{enumerate}
  \item Chiama subito l'operatore per bloccare: China Mobile 10086 / Unicom 10010 / Telecom 10000, chiedi la sospensione immediata e il blocco della sostituzione SIM per impedire all'aggressore di usare il tuo numero.
  \item Accedi agli account critici da un altro dispositivo: da computer o dal telefono di un familiare cambia subito le password di banca, Alipay, WeChat e Apple ID e disconnetti i dispositivi sconosciuti.
  \item Congela i conti di pagamento: Alipay 95188 blocco, Centro sicurezza WeChat blocco account, assistenza banca blocco home banking.
  \item Vai in filiale con il documento: personalmente, con il documento, richiedi una nuova SIM e ripristina l'uso del numero.
  \item Denuncia: il SIM swap configura il reato di «violazione di dati personali dei cittadini»; presenta denuncia con tutte le prove (registrazione della sospensione, perdite economiche, sostituzione SIM).
  \item Controllo completo: dopo aver recuperato il numero, verifica uno per uno gli accessi recenti e i movimenti finanziari per capire se c'è stato uso fraudolento.
\end{enumerate}

\begin{tcbwarn}{Prevenire è meglio che curare — 5 cose quotidiane}
\begin{enumerate}
  \item Password di blocco schermo + autenticazione a due fattori, entrambe attive
  \item Verifica che «Trova il mio telefono» sia attivo
  \item Imposta il PIN della SIM (Impostazioni > Telefono > PIN SIM): se perdi il telefono, nessuno può usare la tua carta
  \item Non installare app di provenienza sconosciuta e non aprire link sospetti
  \item Esegui regolarmente il backup dei dati importanti (foto, rubrica, chat)
\end{enumerate}

\end{tcbwarn}

\begin{tcbpractice}{Pratica: crea la tua «scheda di contatti di emergenza»}
\begin{enumerate}
  \item Prendi un foglio e scrivi questi numeri: assistenza operatore (China Mobile 10086 / Unicom 10010 / Telecom 10000), assistenza banca, linea anti-frode 96110, 110, telefono dei familiari.
  \item Scrivi l'email di registrazione del tuo Apple ID/account del marchio e i contatti di riserva.
  \item Tieni questo foglio nel portafoglio o attaccato al frigorifero — non nel telefono, perché in caso di emergenza il telefono potrebbe non esserci.
  \item Controlla adesso: «Trova il mio telefono» è attivo? Il PIN della SIM è impostato?
\end{enumerate}

\end{tcbpractice}

\begin{tcbsummary}{Riepilogo del capitolo}
\begin{itemize}
  \item Telefono perso, 30 minuti d'oro: chiama da un altro telefono blocca da remoto sospendi SIM congela pagamenti avvisa parenti denuncia
  \item Dopo una truffa, 72 ore: stop trasferimenti conserva prove chiama 96110 denuncia contatta banca cambia password
  \item Telefono infetto: modalità provvisoria disinstalla app sospette scansione ripristino di fabbrica (soluzione estrema)
  \item Account rubato: cambia password disconnetti dispositivi controlla accessi attiva 2FA verifica dati collegati controlla pagamenti
  \item Le estorsioni sono truffe al 99,9\%: non rispondere, non pagare, screenshot e blocca
  \item Truffe con deepfake: osserva battito palpebre/labbra/bordi/luce, chiedi di girare la testa o coprire il viso, fai domande private, richiama un numero noto
  \item SIM swap: se il telefono mostra «Nessun servizio» diffida; imposta PIN SIM, attiva verifica operatore, per gli account importanti usa un'autenticatore
  \item Cinque prevenzioni quotidiane: blocco schermo+2FA, Trova il mio telefono, PIN SIM, niente app sconosciute, backup regolari
\end{itemize}

\end{tcbsummary}
